% ========================================================
% BÀI TẬP TOÁN 5 - KẾT NỐI TRI THỨC VỚI CUỘC SỐNG
% Bài 60: Quãng đường, thời gian của một chuyển động đều
% Định dạng: BTTeX Part III - Trắc nghiệm điền số ngắn
% ========================================================

% Câu 1
\begin{ex}
  Một ô tô đi với vận tốc $50\text{ km/h}$ trong $3$ giờ. Quãng đường ô tô đi được là bao nhiêu ki-lô-mét?
  \shortans{150}
  \loigiai{
    $50 \times 3 = 150\text{ km}$.
  }
\end{ex}

% Câu 2
\begin{ex}
  Một người đi xe máy với vận tốc $40\text{ km/h}$ trong $2$ giờ $30$ phút ($2{,}5$ giờ). Quãng đường đi được là bao nhiêu ki-lô-mét?
  \shortans{100}
  \loigiai{
    $40 \times 2{,}5 = 100\text{ km}$.
  }
\end{ex}

% Câu 3
\begin{ex}
  Một người đi xe đạp với vận tốc $12\text{ km/h}$ trong $45$ phút ($0{,}75$ giờ). Quãng đường đi được là bao nhiêu ki-lô-mét?
  \shortans{9}
  \loigiai{
    $12 \times 0{,}75 = 9\text{ km}$.
  }
\end{ex}

% Câu 4
\begin{ex}
  Một tàu hỏa đi với vận tốc $60\text{ km/h}$ quãng đường $240\text{ km}$. Thời gian tàu hỏa đi hết bao nhiêu giờ?
  \shortans{4}
  \loigiai{
    $240 : 60 = 4\text{ giờ}$.
  }
\end{ex}

% Câu 5
\begin{ex}
  Một người đi xe máy quãng đường $140\text{ km}$ với vận tốc $35\text{ km/h}$. Thời gian đi hết bao nhiêu giờ?
  \shortans{4}
  \loigiai{
    $140 : 35 = 4\text{ giờ}$.
  }
\end{ex}

% Câu 6
\begin{ex}
  Một người đi bộ với vận tốc $5\text{ km/h}$ đi quãng đường $15\text{ km}$. Thời gian đi hết bao nhiêu giờ?
  \shortans{3}
  \loigiai{
    $15 : 5 = 3\text{ giờ}$.
  }
\end{ex}

% Câu 7
\begin{ex}
  Một ca nô đi trên sông với vận tốc $24\text{ km/h}$. Trong $1$ giờ $15$ phút ($1{,}25$ giờ), ca nô đi được bao nhiêu ki-lô-mét?
  \shortans{30}
  \loigiai{
    $24 \times 1{,}25 = 30\text{ km}$.
  }
\end{ex}

% Câu 8
\begin{ex}
  Một máy bay bay với vận tốc $800\text{ km/h}$ quãng đường $2\,000\text{ km}$. Thời gian bay đổi ra phút là bao nhiêu phút?
  \shortans{150}
  \loigiai{
    $2\,000 : 800 = 2{,}5\text{ giờ} = 150$ phút.
  }
\end{ex}

% Câu 9
\begin{ex}
  Một người đi từ $A$ lúc $7\text{ giờ}$ với vận tốc $40\text{ km/h}$ và đi quãng đường dài $80\text{ km}$. Người đó đến nơi lúc mấy giờ?
  \shortans{9}
  \loigiai{
    Thời gian đi: $80 : 40 = 2\text{ giờ}$. Đến nơi lúc $7 + 2 = 9\text{ giờ}$.
  }
\end{ex}

% Câu 10
\begin{ex}
  Một xe khách khởi hành lúc $8\text{ giờ } 15\text{ phút}$ đi quãng đường $120\text{ km}$ với vận tốc $48\text{ km/h}$. Xe đến nơi lúc mấy giờ $45$ phút?
  \shortans{10}
  \loigiai{
    $120 : 48 = 2{,}5\text{ giờ} = 2\text{ giờ } 30\text{ phút}$. Đến nơi lúc $8\text{h}15 + 2\text{h}30 = 10\text{h}45$.
  }
\end{ex}

% Câu 11
\begin{ex}
  Một con ốc sên bò với vận tốc $12\text{ cm/phút}$. Bò trong $15$ phút được bao nhiêu xăng-ti-mét?
  \shortans{180}
  \loigiai{
    $12 \times 15 = 180\text{ cm}$.
  }
\end{ex}

% Câu 12
\begin{ex}
  Một con ong mật bay với vận tốc $8\text{ km/h}$. Trong $15$ phút con ong bay được bao nhiêu ki-lô-mét?
  \shortans{2}
  \loigiai{
    $15\text{ phút} = 0{,}25\text{ giờ}$. $8 \times 0{,}25 = 2\text{ km}$.
  }
\end{ex}

% Câu 13
\begin{ex}
  Quãng đường từ $A$ đến $B$ dài $135\text{ km}$. Một ô tô đi với vận tốc $45\text{ km/h}$, dọc đường nghỉ $30$ phút. Xe khởi hành lúc $6\text{ giờ } 30\text{ phút}$ thì đến $B$ lúc mấy giờ?
  \shortans{10}
  \loigiai{
    Thời gian chạy: $135 : 45 = 3\text{ giờ}$. Tổng thời gian: $3{,}5$ giờ. Đến nơi: $6\text{h}30 + 3\text{h}30 = 10\text{ giờ}$.
  }
\end{ex}

% Câu 14
\begin{ex}
  Một người đi xe đạp từ $A$ đến $B$ với vận tốc $12\text{ km/h}$ hết $2$ giờ $30$ phút. Quãng đường $AB$ dài bao nhiêu ki-lô-mét?
  \shortans{30}
  \loigiai{
    $12 \times 2{,}5 = 30\text{ km}$.
  }
\end{ex}

% Câu 15
\begin{ex}
  Ở câu trên, nếu người đó đi bằng xe máy với vận tốc $30\text{ km/h}$ thì đi hết bao nhiêu giờ?
  \shortans{1}
  \loigiai{
    $30 : 30 = 1\text{ giờ}$.
  }
\end{ex}

% Câu 16
\begin{ex}
  Một vận động viên chạy với vận tốc $5\text{ m/s}$. Trong $1$ phút vận động viên chạy được bao nhiêu mét?
  \shortans{300}
  \loigiai{
    $5 \times 60 = 300\text{ m}$.
  }
\end{ex}

% Câu 17
\begin{ex}
  Một xe lửa đi qua một cây cầu dài $360\text{ m}$ với vận tốc $18\text{ m/s}$. Toa tàu đầu tiên đi từ đầu cầu đến cuối cầu hết bao nhiêu giây?
  \shortans{20}
  \loigiai{
    $360 : 18 = 20\text{ giây}$.
  }
\end{ex}

% Câu 18
\begin{ex}
  Bác Hải đi xe máy từ nhà về quê dài $90\text{ km}$ với vận tốc $45\text{ km/h}$. Bác xuất phát lúc $7\text{ giờ } 15\text{ phút}$. Bác về đến quê lúc mấy giờ $15$ phút?
  \shortans{9}
  \loigiai{
    Thời gian đi: $90 : 45 = 2\text{ giờ}$. Đến lúc $7\text{h}15 + 2\text{h} = 9\text{h}15$.
  }
\end{ex}

% Câu 19
\begin{ex}
  Một ô tô chạy quãng đường $160\text{ km}$ hết $3$ giờ $12$ phút ($3{,}2$ giờ). Vận tốc của ô tô là bao nhiêu km/h?
  \shortans{50}
  \loigiai{
    $160 : 3{,}2 = 50\text{ km/h}$.
  }
\end{ex}

% Câu 20
\begin{ex}
  Một người đi bộ với vận tốc $4\text{ km/h}$. Người đó muốn đi quãng đường $10\text{ km}$ thì hết bao nhiêu phút?
  \shortans{150}
  \loigiai{
    $10 : 4 = 2{,}5\text{ giờ} = 150$ phút.
  }
\end{ex}

% Câu 21
\begin{ex}
  Một chiếc xe đạp đi được $18\text{ km}$ trong $1{,}5$ giờ. Nếu đi trong $2$ giờ thì đi được bao nhiêu ki-lô-mét?
  \shortans{24}
  \loigiai{
    Vận tốc: $18 : 1{,}5 = 12\text{ km/h}$. Quãng đường: $12 \times 2 = 24\text{ km}$.
  }
\end{ex}

% Câu 22
\begin{ex}
  Một người bơi với vận tốc $2\text{ m/s}$. Để bơi hết quãng đường $50\text{ m}$ cần bao nhiêu giây?
  \shortans{25}
  \loigiai{
    $50 : 2 = 25\text{ giây}$.
  }
\end{ex}

% Câu 23
\begin{ex}
  Một ca nô xuôi dòng với vận tốc $32\text{ km/h}$ trong $1$ giờ $30$ phút. Quãng đường ca nô đi được là bao nhiêu ki-lô-mét?
  \shortans{48}
  \loigiai{
    $32 \times 1{,}5 = 48\text{ km}$.
  }
\end{ex}

% Câu 24
\begin{ex}
  Quãng đường dài $175\text{ km}$, xe chạy với vận tốc $70\text{ km/h}$. Thời gian chạy đổi ra phút là bao nhiêu phút?
  \shortans{150}
  \loigiai{
    $175 : 70 = 2{,}5\text{ giờ} = 150$ phút.
  }
\end{ex}

% Câu 25
\begin{ex}
  Một ô tô chở khách đi từ tỉnh $A$ đến tỉnh $B$ dài $210\text{ km}$ hết $3{,}5$ giờ. Vận tốc ô tô là bao nhiêu km/h?
  \shortans{60}
  \loigiai{
    $210 : 3{,}5 = 60\text{ km/h}$.
  }
\end{ex}
